\documentclass[11pt]{article}
\usepackage[a4paper,margin=26mm]{geometry}
\usepackage{amsmath,amssymb,mathtools}
\usepackage{unicode-math}
\setmainfont{Latin Modern Roman}
\setmathfont{Latin Modern Math}
\usepackage{xurl,hyperref}
\hypersetup{hidelinks,pdftitle={Problem statement}}
\newcommand{\sref}[2]{\hyperlink{source:#1:#2}{[S#2]}}
\newcommand{\cataloguescope}[1]{#1}
\begin{document}
\section{Conjecture 2.3 — (X.-S.}\label{statement}

\subsection{Definitions and mathematical statement}

\paragraph{Shared notation.}$\mathbb N=\{1,2,\ldots\}$, and $\mathbb Z,\mathbb Q,\mathbb R,\mathbb C$ denote the integers, rationals, reals and complex numbers. Any other conventions are specified locally.


A knot or link is a smooth oriented embedding of one or more circles in the oriented $3$-sphere, considered up to orientation-preserving ambient isotopy. Write $O$ for the unknot, which bounds an embedded disk. Let $R\mathcal K$ be the free $R$-module on oriented knot isotopy classes. Extend a knot invariant to singular knots with ordinary transverse double points by $v(K_\times)=v(K_+)-v(K_-)$. Let $F_d(R\mathcal K)$ be the submodule generated by the alternating resolution sums of knots with $d$ such points. An $R$-valued finite-type (Vassiliev) invariant of degree at most $d$ is an $R$-linear map $R\mathcal K\to R$ vanishing on $F_{d+1}(R\mathcal K)$. A Jacobi diagram on an oriented circle consists of a finite uni-trivalent graph with univalent vertices attached to the circle and a cyclic ordering of the three incident half-edges at each trivalent vertex. Diagrams are considered up to orientation-preserving identification of the circle and graph. Its degree is half the total number of univalent and trivalent vertices. Let $\mathcal A(S^1;R)$ be the graded $R$-module they generate modulo AS, IHX and STU: reversing a trivalent cyclic order changes sign (AS); the three local reconnections of four external half-edges satisfy the oriented Jacobi identity (IHX); the difference between the two orders of neighboring legs on the circle equals the diagram joining them at a trivalent vertex (STU). All relations keep the rest of the diagram fixed. The framing-independence relation FI additionally sets to zero a diagram with an isolated chord whose two endpoints are consecutive on the circle. Write $\mathcal A_d$ for degree $d$. The canonical resolution map $\eta_d:\mathcal A_d(S^1;R)/\mathrm{FI}\to F_d(R\mathcal K)/F_{d+1}(R\mathcal K)$ sends a diagram to its corresponding singular-knot resolution class; trivalent diagrams are interpreted using STU. The weight system of an invariant $v$ of degree at most $d$ is the composition of $\eta_d$ with the map induced by $v$ on this graded piece.
\paragraph{Statement recorded in the supplied source.}
For every commutative ring $R$ with identity, every integer $d\ge0$, and every $R$-linear weight system
\[w:\mathcal A_d(S^1;R)/\mathrm{FI}\longrightarrow R,\]
is there a degree-at-most-$d$ Vassiliev invariant $v:R\mathcal K\to R$ whose degree-$d$ weight system is $w$? Include positive-characteristic coefficient rings; the question is not restricted to $R=\mathbb Q$. \sref{statement}{2}
\subsection{Short English statement}
Does every weight system over every commutative coefficient ring integrate to a finite-type knot invariant?
\subsection{Sources}
\begin{description}
\item[\hypertarget{source:statement:1}{S1}] ULAM. AI and the UnsolvedMath Contributors, \emph{UnsolvedMath: A Curated Collection of Open Mathematics Problems}, catalogue entry. CC BY 4.0. \url{https://huggingface.co/datasets/ulamai/UnsolvedMath}.
\item[\hypertarget{source:statement:2}{S2}] T. Ohtsuki (editor), \emph{Problems on invariants of knots and 3-manifolds}, Geometry \& Topology Monographs 4 (2002), pp. 377--572 (published 2004), item 2.3 and its surrounding definitions and remarks. \url{https://arxiv.org/abs/math/0406190v1}.
\end{description}
\label{end:statement}
\end{document}
